\chapter{Zur kirchlichen Chronik}
\quellseite{6}

\section{America}

Wie sehr der Hausgottesdienst in den christlichen Familien abgenommen hat, würde eine genaue Nachfrage in den einzelnen Häusern bald herausstellen. In allen Kirchengemeinschaften wird darüber geklagt. Und doch ist solcher Hausgottesdienst so überaus nöthig und wichtig für die Wohlfahrt der Kirche und der Familie. Die Seele des Christen erhält nicht ihre Gebühr, wenn sie nicht täglich aus Gottes Wort unterrichtet, ermahnt, ermuntert und getröstet wird. Sie verkümmert, weil ihr die Luft fehlt. Und wie nothwendig ist besonders dem heranwachsenden jüngeren Geschlecht der tägliche Umgang mit Gottes Wort und Gebet zur Stärkung wider die mancherlei Gefahren und Versuchungen in der Welt. Werden sie darin eifrig und fleißig sein, wenn sie sehen, daß ihre Eltern sich begnügen lassen an dem oft nicht einmal regelmäßigen Besuch des öffentlichen Gottesdienstes? An der zum Hausgottesdienst erforderlichen Zeit fehlt es in keinem Hause, wohl aber in vielen an dem Wollen. Gottes Wort ermahnt: „Lasset das Wort Christi unter euch reichlich wohnen in aller Weisheit. Lehret und vermahnet euch selbst mit Psalmen und Lobgesängen und geistlichen lieblichen Liedern, und singet dem Herrn in eurem Herzen“, Col. 3, 16. Gott gebe dir dazu, lieber Leser, Wollen und Vollbringen. L. F.

New Orleans, La. Die letzte Nummer der „Evangelisch-Lutherischen Blätter“ bringt den Jahresbericht aus unsern New Orleanser Gemeinden. Hiernach bestehen in New Orleans und den Vorstädten, die drei Negergemeinden eingeschlossen, zehn lutherische Gemeinden. Die Zahl der stimmberechtigten Glieder ist verhältnißmäßig nicht groß. Sie beträgt nur 359. Dies hat seinen Grund in den eigenthümlichen Verhältnissen. In New Orleans, wie im Süden überhaupt, entfaltet das verderbliche Logenwesen eine ganz besondere Macht und hält namentlich die Männer in der Allerweltsreligion, die keine Seligkeit durch Christum den Gekreuzigten will, gefangen. Dennoch ist die Zahl der lutherischen Christen, die in der Stadt New Orleans mit uns denselben Glauben bekennen, durch Gottes Gnade nicht so gering. Zu den Gemeinden gehören 3766 communicirende Glieder. Die Zahl der Seelen beträgt 6008. In den Gemeindeschulen sind 877 Kinder. Confirmirt wurden (10 Erwachsene eingeschlossen) 244. Auch die 1735 Kinder, welche die Sonntagsschulen besuchten, werden aus Gottes Wort gelernt haben, wie man recht glauben, christlich leben und selig sterben kann. Es ist ja traurig und will uns entmuthigen, wenn wir an einigen Orten in unserer Synode die Erfahrung machen, daß die Gemeinden klein bleiben an stimmberechtigten Gliedern. Aber vergessen wir nicht, daß der Befehl Christi nicht zunächst dahin lautet: „Bildet große, wohlgeordnete Gemeinden“, sondern: „Prediget das Evangelium.“ Wo wir Gelegenheit haben, diesen Befehl auszurichten, sollen wir mit großer Freudigkeit arbeiten und alles andere Gott befehlen. F. B.

In Bakersfield, Cal., hat sich, wie wir aus dem „Lutherischen Botschafter“ ersehen, Anfangs dieses Jahres eine lutherische Gemeinde organisirt. Die Gemeinde ist freilich vorläufig klein. Sie zählt nur 13 stimmberechtigte Glieder. Aber Bakersfield ist ein Eisenbahnknotenpunkt im südlichen California, wo eine bedeutende Anzahl Arbeiter Beschäftigung findet. Gott gebe, daß von hier aus das Wort des Evangeliums noch weiter in die Lande leuchte. F. B.

Ueber die Mission unter den zerstreut wohnenden lutherischen Lithauern, die hauptsächlich innerhalb unsers Oestlichen Districts seit August von P. Brukst in Waterbury, Conn., betrieben wird, findet sich ein Bericht in einer der letzten Nummern des „Zeuge und Anzeiger“, aus dem wir Folgendes herausheben. Der Missionar schreibt: „Wie bei allen Missionen, so gibt es auch bei der lithauischen Mission Schwierigkeiten aller Art, was ja gar nicht anders erwartet werden kann. Im Allgemeinen muß man wohl sagen, daß viele unter den Lithauern sind, die kein Verlangen nach dem Worte Gottes und nach ihrer Seelen Seligkeit haben. Sie sind kalt, weil sie in diesem Lande so lange ohne Gottes Wort dahingegangen sind. Aber die meisten haben Hunger und Durst nach dem Brod und Wasser des Lebens und sind um ihr Seelenheil besorgt. Ja, man findet oft Leute, die mit Thränen in den Augen ihr großes Verlangen nach dem Wort Gottes, nach dem süßen Evangelium zeigen; so daß man Mühe hat, sie zu beruhigen und zu trösten. — Schwer ist bei dieser Mission, daß die Leute so weit zerstreut wohnen. Denn die vielen großen und oft umständlichen Reisen nehmen nicht nur meine meiste Zeit in Anspruch, sondern sind auch mit großen Unkosten verbunden. Und die Lithauer, obwohl sie sich anstrengen, können doch nicht viel für meine Reisen aufbringen, da sie meist sehr arm sind und große Familien zu ernähren haben. — Bis jetzt habe ich neun Missionsstationen, nämlich: Waterbury, Naugatuck und Bridgeport, Conn.; Boston, Mass.; Brooklyn, N. Y.; Scranton, Wilkesbarre und Pittsburg, Pa.; Cleveland, O. Die Lithauer an diesen Orten haben mich schriftlich oder mündlich als ihren Pastor und Seelsorger berufen und wünschen, daß ich sie so oft als möglich bediene. Es gibt noch an andern Plätzen Lithauer, besonders in Pennsylvania, die ich aber noch nicht besucht habe. Die Leute bitten mich dringend, auch ihnen das Brod des Lebens zu bringen; aber woher soll man die Zeit und die Mittel dazu nehmen? Um Zeit und Reisekosten zu sparen, habe ich meinen Arbeitsplan so eingerichtet, daß ich einen Sonntag in Waterbury und Naugatuck, Conn., Gottesdienst halte und den zweiten Sonntag verreise, so daß die näher liegenden Stationen öfters und die ferner liegenden etwas seltener bedient werden. Ich habe nun von August bis December schon neun zum Theil große und lange Reisen gemacht, die sich bis nach Scranton, Wilkesbarre, Pittsburg, Pa., und Cleveland, O., ausdehnten; habe
\quellseite{7}
siebenmal Beichtgottesdienst gehalten und 91 Personen das heilige Abendmahl gereicht, sowie andere Amtshandlungen verrichtet. In Pawtucket werden 10 Kinder jeden Mittwoch-Nachmittag von vier Uhr an unterrichtet, und Abends ist Confirmandenunterricht. Von Waterbury kann ich diesmal nur berichten, daß unser Werk sehr langsam vorangeht. Der Hauptgrund hiervon ist, daß unser Gottesdienstlocal den Leuten zu entlegen ist, sie haben alle nicht weniger als vierzig bis fünfzig Minuten zu gehen. Dazu kommt noch, daß wir unsere Gottesdienste nur Nachmittags halten können, da die Kirche Vormittags anderweitig im Gebrauch ist; auch das hält viele zurück. Ich habe mich bemüht, ein anderes Local zu bekommen, aber vergeblich. Bis jetzt ist es mir auch nicht gelungen, ein Zimmer zu finden, wo ich die Kinder wenigstens einige Abende in der Woche unterrichten könnte.“ — Gott segne auch diese Mission und führe durch sie viele lithauische Glaubensgenossen zum ewigen Leben! L. F.

In der Episkopalkirche unseres Landes wird die Hinneigung zum Papstthum und seinen antichristlichen Greueln immer allgemeiner, und die sogenannte ritualistische oder hochkirchliche Richtung gewinnt immer mehr Boden. Allein in Chicago gibt es neun Kirchen dieser Gemeinschaft, in denen täglich das Meßopfer dargebracht wird. Vor zwanzig Jahren noch fanden sich kaum zwei oder drei Episkopalkirchen in unserm Lande, die eine Vorrichtung für die Ohrenbeichte hatten, jetzt sind es ihrer schon mehr als dreißig. Viele Glieder der Episkopalkirche beten regelmäßig den Rosenkranz, und gar manche unterscheiden sich nur noch durch den Namen von den Anhängern der Papstkirche. Auch diese Thatsache muß bedacht werden, wenn man die Macht und den Einfluß der römischen Kirche in unserm Lande ins Auge faßt. Sie hat viele Gesinnungsgenossen in einer andern Kirche, die gleichfalls besonders reich an irdischem Vermögen und Einfluß ist. L. F.

\subsection{Ein schweres Hinderniß für die protestantische Mission auf den Philippinen}

ist die große Gleichgültigkeit der dort wohnenden Americaner gegen das gepredigte Wort. Ein protestantischer Zeitungsschreiber, der sich zwei Jahre lang auf den Inseln aufgehalten und ganz besonders auch die religiösen und kirchlichen Zustände untersucht hat, theilt mit, daß von den 8000 gegenwärtig in Manila sich aufhaltenden Americanern nur ungefähr 50 Sonntags die Kirche besuchen, während die katholischen Filipinos gewohnt sind, ganz regelmäßig an der römischen Messe Theil zu nehmen. Diese schließen dann, daß das, was von den Protestanten selbst so gering geschätzt wird, nichts Besseres sein könne als ihre bisherige Religion, die freilich nichts anderes als ein etwas christlich aufgeputztes Heidenthum ist. Die Folge davon ist, daß sie von protestantischen Missionaren wenig wissen wollen und in ihrem römisch-heidnischen Aberglauben verharren. Auch die eigentliche Heidenmission wird durch den ungöttlichen Wandel derer, die den Christennamen tragen, sehr gehemmt. So erfüllt sich auch dort das Schriftwort: „Eurethalben wird Gottes Name gelästert unter den Heiden“, Röm. 2, 24. L. F.

\section{Ausland}

Zu Lebzeiten des vorigen Papstes, Leos XIII., hatten sich süditalienische Priester an den Papst mit einem Bittschreiben gewandt, er möchte doch das Verbot der Ehelosigkeit für den Priesterstand der römischen Kirche aufheben. Dieses Bittgesuch war damals dem von Jesuiten umringten und geleiteten Leo gar nicht vorgelegt worden. Vor Kurzem wurde nun berichtet, daß dieselben Bittsteller mit erneuter Hoffnung den neuen Papst, Pius X., angehen wollten. Ihre Hoffnung wird aber ohne Zweifel wieder zu Schanden werden. Denn Pius hat schon öffentlich erklärt und auch deutlich gezeigt, daß er ganz in den Bahnen seiner Vorgänger wandeln will. Und doch wäre, gerade auch in diesem Stücke, wie selbst Katholiken zugeben, eine Reform dringend nötig. Die „Wartburg“ theilt mit, daß erst in jüngster Zeit der Osservatore Cattolico, das Organ des römischen Cardinals Rampolla, in einem Aufsatz eine eingehende Betrachtung widmete „den überaus schmerzlichen Vorkommnissen, daß innerhalb weniger Tage zwei Geistliche durch Selbstmord endigten, daß mehrere Mädchen grausame Rache mit bewaffneter Hand an geistlichen Verführern nahmen, daß zahlreiche Richtersprüche zahlreiche Verurtheilte als unwürdig des Priestergewandes kennzeichneten, ja, daß verschiedene zur Verdeckung von Sittlichkeitsverbrechen verdächtige Morde von Pfarrern zu beklagen seien“. Mit Schärfe wandte sich der Osservatore Cattolico gegen die Selbsttäuschung einiger katholischen Blätter, durch Todtschweigen die Uebelthaten verdecken zu wollen. — Die Schrift bezeichnet es als eine Teufelslehre, zu verbieten, ehelich zu werden, 1 Tim. 4, 1. 3., und unser lutherisches Bekenntniß nennt darum mit Recht das päpstliche Eheverbot ein Kennzeichen des Antichrists, sagt: „Darum ist das Gesetz, dadurch die Priesterehe verboten, ein recht Papstgesetz der römischen Tyrannei“ und weist nach, wie diese Tyrannei zu viel Unzucht und Ehebruch, zu allerlei Schanden und Lastern führt: (Apologie, Müller, S. 240. 245.) Das hat die Geschichte des Papstthums je und je bestätigt bis auf den heutigen Tag. L. F.

Wie groß der Aberglaube auch in der griechisch-katholischen Kirche sich äußert, zeigt, was neuerdings in der „Allgemeinen Ev.-Luth. Kirchenzeitung“ aus Rußland berichtet wird. Die russische Kirche hat vor einiger Zeit einen neuen Heiligen ernannt, Serafim. In Folge dessen strömen die Pilger an die Grabstätte des „Heiligen“ in Sarowo, und die dortige Kathedrale ist täglich von ihnen gefüllt. Täglich werden auch nach den verschiedensten Theilen des russischen Reiches nicht nur Bilder des Heiligen, sondern auch Wasser aus der „heiligen“ Quelle und sogar Pakete „heiligen“ Sandes versandt, so daß die Post in Sarowo kaum allen an sie gestellten Anforderungen nachkommen kann. Ja, manche Pilger haben sogar den Versuch gemacht, von dem Sarge, in welchem der „Heilige“ gebettet ist, kleine Stückchen abzubeißen, so daß der Sarg aus Vorsicht mit einem Metallüberzuge versehen werden mußte. — Armes, betrogenes Volk, das die Heiligen um Hülfe anruft und damit eine zweifache Sünde thut: den Herrn, die lebendige Quelle, verlassen sie, und machen sich hie und da ausgehauene Brunnen, die doch löchericht sind und kein Wasser geben, Jer. 2, 13. L. F.

\section{Aus Welt und Zeit}

„Welches ist der wirkliche Himmel für den Arbeitsmann?“ fragte vor einiger Zeit der bekannte englische Staatsmann und Schriftsteller John Morley in einer Rede, die er zu Brechin in Schottland vor einer zahlreichen Zuhörerschaft hielt. Er antwortete dann selbst und sagte: „Es ist ständige Arbeit, für die man einen guten Lohn bekommt“, und mit stürmischem Beifall wurde diese Beschreibung des „wahren Himmels“ von der Menge begrüßt. Ganz ähnlich äußerte sich vor Kurzem nach dem Bericht des „Lutherischen Kirchenboten für Australien“ der Präsident des Handwerkervereins in Melbourne, Australien, und sagte vor einer großen Versammlung: „Meine Religion ist Arbeit; ich weiß von keiner andern Religion.“ Aehnlich denken und reden Hunderttausende in unserer Zeit. Das ist, wie Arbeiterführer gesagt haben, „das Evangelium des neunzehnten Jahrhunderts“, und überall in der Welt hören wir die Losung der Arbeiter:

\begin{verse}
Wir wollen auf Erden glücklich sein\\
Und wollen nicht mehr darben.
\end{verse}

Aber das ist eine elende Religion, ein erbärmliches Evangelium, ein trauriger Himmel. Christlicher Arbeiter Evangelium ist das
\quellseite{8}
uralte und doch ewig neue \gesperrt{Evangelium} von Christo, das „eine Kraft Gottes ist, die da selig macht alle, die daran glauben“, Röm. 1, 16.; ihre \gesperrt{Religion} ist: „Ich lebe; doch nun nicht ich, sondern Christus lebet in mir. Denn was ich jetzt lebe im Fleisch, das lebe ich in dem Glauben des Sohns Gottes, der mich geliebet hat und sich selbst für mich dargegeben“, Gal. 2, 20.; und ihr \gesperrt{Himmel} ist der neue Himmel und die neue Erde, „in welchen Gerechtigkeit wohnet“, und wo Gottes Volk „wohnen wird in Häusern des Friedens, in sichern Wohnungen und in stolzer Ruhe“, 2 Petr. 3, 13. Jes. 32, 18. Dabei erhalte Gott alle christlichen Arbeiter und bewahre sie vor den genannten Zügen unserer Tage und vor dem Schicksal, das denen bevorsteht, die da lieb haben und thun diese Züge. „Wie es geschah zu den Zeiten Noah, so wird’s auch geschehen in den Tagen des Menschensohns. Sie aßen, sie tranken, sie freieten, sie ließen sich freien bis auf den Tag, da Noah in die Arche ging, und kam die Sintfluth und brachte sie alle um. Deselbigen gleichen, wie es geschah zu den Zeiten Lot. Sie aßen, sie tranken, sie kauften, sie verkauften, sie pflanzten, sie baueten. An dem Tage aber, da Lot aus Sodom ging, da regnete es Feuer und Schwefel vom Himmel, und brachte sie alle um. Auf diese Weise wird’s auch gehen an dem Tage, wenn des Menschen Sohn soll offenbar werden“, Luc. 17, 26. ff. L. F.
